% Table 1: simulation evidence for pooled estimation and uncertainty assessment.
\begin{table}[!p]
\centering
\fontsize{8.5}{9.5}\selectfont
\setlength{\smallskipamount}{0.8pt}
\caption{Simulation evidence for pooled estimation and uncertainty assessment.
(A) Coefficient-specific root mean squared error (RMSE), with signed relative
bias (\%) in parentheses. (B) Mean SE / empirical SD and 95\% coverage for
three covariance calculations. (C) Bias, RMSE, and 95\% coverage for four
interval procedures on matched datasets. B--C report minima--maxima across
six regression coefficients, including the intercept.}
\label{tab:sim-essential}

\smallskip
\begin{minipage}{\textwidth}
\small
\textbf{A. Estimation from short episodes: RMSE and relative bias (\%)}\par
$N=150$, $T=30$, $\sigma_\eta^2=1$; $M=60$ per noise level.
State-space fits also estimate $\sigma_\epsilon^2$.
\end{minipage}\par\smallskip
{\setlength{\tabcolsep}{2.4pt}
\renewcommand{\arraystretch}{0.90}
\begin{tabular*}{\textwidth}{@{\extracolsep{\fill}}l*{9}{r}@{}}
\toprule
& \multicolumn{3}{c}{$\sigma_\epsilon^2=0.25$}
& \multicolumn{3}{c}{$\sigma_\epsilon^2=1$}
& \multicolumn{3}{c}{$\sigma_\epsilon^2=4$}\\
\cmidrule(lr){2-4}\cmidrule(lr){5-7}\cmidrule(l){8-10}
Coefficient
& Joint & \shortstack{OLS\\two-stage} & \shortstack{State-space\\two-stage}
& Joint & \shortstack{OLS\\two-stage} & \shortstack{State-space\\two-stage}
& Joint & \shortstack{OLS\\two-stage} & \shortstack{State-space\\two-stage}\\
\midrule
@@PANEL_A@@
\bottomrule
\end{tabular*}}

\smallskip
\begin{minipage}{\textwidth}
\small
\textbf{B. Contribution of observed information: SE accuracy and coverage}\par
$N=1174$, $G=44$, $T=14$; $M=500$ per $\tau$;
$\sigma_\epsilon^2=400$ in generation; all ten parameters estimated.
\end{minipage}\par\smallskip
{\setlength{\tabcolsep}{2.5pt}
\renewcommand{\arraystretch}{0.90}
\begin{tabular*}{\textwidth}{@{\extracolsep{\fill}}llrrrr@{}}
\toprule
& & \multicolumn{2}{c}{Model-based} & \multicolumn{2}{c}{CR0}\\
\cmidrule(lr){3-4}\cmidrule(l){5-6}
$\tau$ & Information calculation & SE / emp.\ SD & Coverage & SE / emp.\ SD & Coverage\\
\midrule
@@PANEL_B@@
\bottomrule
\end{tabular*}}

\smallskip
\begin{minipage}{\textwidth}
\small
\textbf{C. Participant-clustered estimation and inference}\par
Same $N,G,T$ as B; $B=999$; observation variance fixed at zero or estimated when positive.
\end{minipage}\par\smallskip
{\setlength{\tabcolsep}{1.6pt}
\renewcommand{\arraystretch}{0.90}
\begin{tabular*}{\textwidth}{@{\extracolsep{\fill}}cccrrrrrr@{}}
\toprule
& & & \multicolumn{2}{c}{Estimation} & \multicolumn{4}{c}{Coverage}\\
\cmidrule(lr){4-5}\cmidrule(l){6-9}
\shortstack{Generating\\$\sigma_\epsilon^2$} & $\tau$ & $M$ & Bias & RMSE & Model-based & CR0 & MD-style & Bootstrap\\
\midrule
@@PANEL_C@@
\bottomrule
\end{tabular*}}

\smallskip
\begin{minipage}{\textwidth}
\footnotesize
\textit{Notes.}
$M$: simulated datasets per setting; $B$: participant-bootstrap resamples per dataset.
In A, $\boldsymbol{\theta}_\phi=(0.30,0.15,-0.20,0.25)^\top$ in row order.
Joint denotes pooled estimation. Two-stage methods regress transformed episode
persistence estimates on covariates after separate OLS fits (ignoring observation
error) or state-space likelihood fits. Relative bias is
$100(\overline{\hat\theta}_{\phi,j}-\theta_{\phi,j})/\theta_{\phi,j}$,
where the bar averages $M$ estimates.
For B--C, coefficient order is intercept, age, sex, BMI, prediabetes, and diabetes.
Generating parameters are
$\boldsymbol{\theta}_\phi=(1.551,0.045,0.092,0.018,0.096,0.336)^\top$
and $(a_1,P_1,\sigma_\eta^2,\sigma_\epsilon^2)=(53.4,1910,38.5,400)$
for B and the positive-variance rows of C, and
$\boldsymbol{\theta}_\phi=(1.641,0.032,0.088,0.037,0.051,0.250)^\top$
and $(a_1,P_1,\sigma_\eta^2,\sigma_\epsilon^2)=(56.5,1739,33.8,0)$
for the zero-variance rows of C.
$\tau$ is the participant random-intercept SD on the persistence link.
In B, SE ratios are mean estimated SE / empirical SD; ignoring nuisance uncertainty
treats $(a_1,P_1,\sigma_\eta^2,\sigma_\epsilon^2)$ as fixed only in covariance calculations.
In C, bias (on the coefficient scale) and RMSE describe the pooled estimates shared
by all four interval procedures. Bias, RMSE and coverage use generating coefficients
at $\tau=0$ and each specification's working-likelihood target at $\tau=0.4$.
Intervals use normal quantiles (model-based), $t_{43}$ (CR0 and MD-style),
or bootstrap percentiles.
\end{minipage}
\end{table}
